\chapter*{}
\begin{center}
    {\Large{\textbf{STRUCTURAL AND HYDROLOGICAL DAM MONITORING SYSTEM USING IOT AND COMPUTER VISION ON EDGE COMPUTING}}}
\end{center}

\vspace{1cm}

\begin{flushright}
    Author: Gladerson Pereira Belizio da Silva\\
\end{flushright}

\vspace{1cm}

\begin{center}
    \Large{\textsc{\textbf{Abstract}}}
\end{center}

\noindent 
The management of water resources and the guarantee of dam safety are critical aspects for the sustainability of the Brazilian semiarid region. In this context, this work proposes a cost-effective architecture, based on the Internet of Things (IoT) and Artificial Intelligence (IA), applied to the monitoring of the Oiticica Dam (RN). The system integrates level sensors based on radar technology and image processing at the edge (Edge Computing), employing single-board microcomputers and computer vision techniques based on deep learning for the autonomous detection of anomalies in concrete structures. Data visualization is consolidated in a supervision platform that correlates real-time telemetry to a georeferenced three-dimensional representation of the monitored infrastructure.

\noindent\textit{Keywords}: IoT; Computer Vision; Edge Computing; Dam Monitoring; Deep Learning.
